\documentclass[10pt,a4paper]{amsart}
\usepackage{fontspec,amsmath,amssymb,amsthm,mathtools,stmaryrd,physics,braket,thmtools,graphicx}
\usepackage[dvipsnames,cmyk]{xcolor}
\usepackage[margin=22mm]{geometry}
\usepackage{youngtab}
\usepackage[numbers]{natbib}
\usepackage{hyperref}
\usepackage{cleveref}
\let\autocite\cite\let\textcite\citet
\declaretheoremstyle[
    spaceabove=\topsep,
    spacebelow=\topsep,
    headindent=0pt,
    headfont=\bfseries,
    postheadhook = {\hspace*{\parindent}},
    postheadspace=0pt plus 1pt minus 1pt,
    headpunct=:,
    notefont=\bfseries,
]{theorem}

\declaretheorem[style=theorem, numberwithin=section, name=Definition]{defn}
\declaretheorem[style=theorem, numberlike=defn, name=Example]{ex}
\declaretheorem[style=theorem, numberlike=defn, name=Theorem]{thrm}
\declaretheorem[style=theorem, numberlike=defn, name=Proposition]{prop}
\declaretheorem[style=theorem, numberlike=defn, name=Lemma]{lem}
\declaretheorem[style=theorem, numberlike=defn, name=Corollary]{cor}
\declaretheorem[style=theorem, numberlike=defn, name=Conjecture]{conj}
\declaretheorem[style=theorem, numberlike=defn, name=Question]{question}
\declaretheorem[style=theorem, numbered=no, name=Remark]{remark}

\declaretheoremstyle[
    spaceabove=\topsep,
    spacebelow=12pt,
    headindent=0pt,
    headfont=\bfseries,
    postheadhook = {\hspace*{\parindent}},
    postheadspace=0pt plus 1pt minus 1pt,
    headpunct=:,
    notefont=\bfseries,
    qed=\qedsymbol,
]{solution}
\declaretheorem[style=solution, numbered=no, name=Proof]{pf}

% change footnote style to symbols instead of numbers
\renewcommand{\thefootnote}{\fnsymbol{footnote}}

% -------------------------------------------------------------------------
%  Common abbreviations and words with accents
% -------------------------------------------------------------------------

% ---- LATIN ----
\newcommand*\ie{i.e.}
\newcommand*\cf{c.f.}
\newcommand*\eg{e.g.}
\newcommand*\apriori{a~priori}

% ---- PROPER NAMES (because I'm lazy) ----
\newcommand*\Burgisser{B\"urgisser}

% -------------------------------------------------------------------------
%  Math operators and symbols
% -------------------------------------------------------------------------

% ---- Paper Specific ----
\DeclareMathOperator{\Sub}{Sub}
\DeclareMathOperator{\pfaff}{pfaff}
\DeclareMathOperator{\sym}{sym}
\DeclareMathOperator{\lex}{lex}
\DeclareMathOperator{\leqlex}{\leq_{\lex}}
\DeclareMathOperator{\geqlex}{\geq_{\lex}}
\DeclareMathOperator{\characteristic}{char}
\newcommand{\svdots}{\raisebox{-0.12\height}[0pt][0pt]{$\vdots$}}
\apptocmd{\det}{\nolimits}{}{}

% ---- Misc Functions ----
\DeclareMathOperator{\sgn}{sgn}

% ---- Matrices ----
\DeclareMathOperator{\IMM}{IMM}

% ---- Polynomials ----
\DeclareMathOperator{\poly}{poly}
\DeclareMathOperator{\coeff}{coeff}

% ---- Sets ----
\newcommand*\N{\mathbb{N}}
\newcommand*\Z{\mathbb{Z}}
\newcommand*\R{\mathbb{R}}
\newcommand*\C{\mathbb{C}}
\newcommand*\F{\mathbb{F}}

% ---- Symbols ----
\newcommand*\e{\varepsilon}

% ---- RELATORS & ARROWS ----
\newcommand*\defeq{\coloneqq}

% -------------------------------------------------------------------------
% Enclosures
% -------------------------------------------------------------------------

% ---- FUNCTIONS ----
\newcommand*\floor[1]{%
    \left\lfloor #1 \right\rfloor
}

% ---- COMPLEXITY MACROS ----
\newcommand*\bigO[1]{%                       % Constants no longer matter
    O\kern-0.0pt\left(#1\right)
}

\newcommand*\bigTheta[1]{%
    \Theta\kern-0.0pt\left(#1\right)
}

\usepackage{float}
